\documentclass[12pt,a4paper]{ctexart}
\usepackage{amsmath,amssymb}
\usepackage{braket}
\usepackage{geometry}
\geometry{margin=2.5cm}

\title{\textbf{人际理解的一种量子力学类比}\\——算符、表象与本征值}
\author{张明楷\thanks{厦门大学物理系。}\quad 语晴\thanks{共一。}}
\date{2026年6月}

\begin{document}

\maketitle

\begin{abstract}
本文尝试借用量子力学的数学结构——算符、本征值、幺正变换与矩阵元——构建一套理解人际差异与包容的类比框架。核心论点是：每个人对应于一个独特的算符，同一外部事件作用其上，将坍缩出不同的本征值。理解他人的实质，不是修改其矩阵元，而是在自身表象下进行幺正变换以求得其本征值，并接受矩阵元不可通约的事实。
\end{abstract}

\section{基本设定：每个人是一个算符}

设每个人 $P$ 对应于一个线性算符 $\hat{P}$，作用在某个公共的状态空间 $\mathcal{H}$ 上。该算符由其\emph{矩阵元}构成——成长环境、创伤经验、被爱的方式、读过的书、走过的路。这些矩阵元在 $\hat{P}$ 被构造之后即部分固定、部分可变。

同一外部事件 $\ket{\psi}$（例如分别、冲突、选择）作用在不同算符上，将产生不同本征值：
\begin{equation}
\hat{P}_A \ket{\psi} = \lambda_A \ket{\phi_A},\qquad
\hat{P}_B \ket{\psi} = \lambda_B \ket{\phi_B}.
\end{equation}
其中 $\lambda_A \neq \lambda_B$ 是常态，并不蕴含任何一方有误。两者满足相同的边界条件——例如`继续向前生活''——即可视为有效的解。

\section{本征值而非矩阵元}

两个人的差异体现在两个层面：

\begin{enumerate}
\item \textbf{本征值层面}：对同一事件的不同解释或行为输出。
\item \textbf{矩阵元层面}：其世界观中具体的、细微的、不可通约的纹理。
\end{enumerate}

理解另一个人的目标，并非要求其矩阵元与我的一致——这是不可能的。目标是求得其\emph{本征值}：在其自身的内部逻辑下，为何选择了这一条路，为何做出了这一个反应。

求得本征值所需的，不是修改对方，而是在\emph{自身的表象}下进行推理：
\begin{equation}
\text{我看到的对方的选择} \;\longrightarrow\; \text{我推测对方选择的原因} \;\longrightarrow\; \text{我发现此原因在其世界观下是自洽的}.
\end{equation}

\section{幺正变换与表象}

幺正变换保持本征值不变，但不保持矩阵元不变。将算符 $\hat{P}_B$ 幺正变换到 $\hat{P}_A$ 的表象下：
\begin{equation}
\hat{P}'_B = U^\dagger \hat{P}_B U,
\end{equation}
得到的是 $\hat{P}_B$ 在 $A$ 的语言中的一种翻译。本征值 $\lambda_B$ 在变换下保持不变：\textbf{对方的核心逻辑，只要我花足够的心力去翻译，是可以理解的。}但变换后的矩阵元 $\bra{i}\hat{P}'_B\ket{j}$ 与原矩阵元 $\bra{i}\hat{P}_B\ket{j}$ 不同：\textbf{你在对方的坐标系里看到的世界，永远只是你的翻译，不是对方的原始视觉。}

这里有一个容易混淆的细微之处：幺正变换应当在\emph{自身表象}下进行，而非在\emph{对方表象}下。因为对方的表象——其童年、其每一次心跳加速与每一次决定沉默的瞬间——是我永远无法真正进入的。那是对方算符的本征表象，天然封闭于我的观测之外。我所能做的，是接受这一封闭性：不试图钻入对方的坐标系，而是在自己的语言框架里完成翻译。物理学不会要求一个观测者变成被观测的粒子；它只要求观测者用自己的仪器，逼近测量值。

试图进入他人表象的行为——在人际中表现为强行闯入对方的心理边界——不是幺正变换，而是对算符自身的扰动。扰动越大，测得的本征值越失真。真正的理解，是站在自己的坐标系里，承认变换的局限性，然后最大程度地逼近对方的本征谱。

这就是为什么人际理解有一个不可逾越的边界：

\begin{itemize}
\item 本征值可以求得：对方为什么这样做；
\item 矩阵元不可知：对方具体的感觉、每一帧画面的质感、那些说不清的细节。
\end{itemize}

理解不等于全知。理解一个人，是找到其本征值的近似，并承认矩阵元永远封闭在那个人的内部。

这里隐含着一个更深层的代价：幺正变换保本征值，却牺牲了矩阵元的结构特征。本征值告诉我她\emph{做了什么选择}——拉黑、沉默、走向新生活。但矩阵元告诉我她\emph{如何体验这些选择}——拉黑时指尖有没有犹豫、沉默里压着愤怒还是释然、走向新生活时是轻盈的还是沉重的。这些纹理在变换之下消散了。我得到的是一张正确的谱，却失去了一支曲子的音色。

这是所有理解的宿命，而非失败的标志。物理学里，两个不对易的算符无法同时对角化。人际中，两个人的表象无法完美重合。求得了本征值，便已走到了变换所能抵达的极限。再往前，就是对他人的暴力——强行要求对方在你的坐标系里交出她的矩阵元。那不是理解，是入侵。

从数学本质上讲，试图由本征值反推他人矩阵元的形成原因，是一个不可行的逆问题。一组本征值并不唯一确定一个算符：无穷多个不同的矩阵可以共享同一组本征值，但矩阵元的结构——即那些构成对方之所以为对方的具体纹理——却迥然不同。同样的\emph{拉黑}行为（本征值），可以来自自我保护、可以来自厌倦、可以来自新男友的要求、可以来自某个深夜突然涌上的释然。本征值告诉你\emph{是什么}，却不告诉你\emph{为什么}。要知道为什么，需要矩阵元。而矩阵元在你变换到自身表象的那一刻，已经不可逆地改变了。

这就是为什么\emph{理解}永远不等于\emph{全知}。你能算出她做了什么，但你算不出她为什么那么做——至少不能以数学意义上的精确性。你只能逼近、猜测、用你自己的经验外推。逼近到一定程度，便是理解。再追，便是徒劳。

\section{刚性矩阵元与柔性矩阵元}

某些矩阵元一旦形成便不再变化。创伤、深层恐惧、边界感——这些是刚性的。它们构成了对方地形的一部分。行走其间，不能期望地形为你改变；你只能绕行。

另一些矩阵元——例如自我认知、对过去的解释框架、与自我的和解程度——则可以是柔性的。它们在你自身的持续反思中逐步软化，表现为同一个人在不同的时间点，对同一事件产生了新的理解。

最好的包容，不是改变他人的刚性边界，而是识别其刚性边界并承认它的真实。同时，不断修养自己的柔性矩阵元，使其在遭遇他人的不同输出时，不立即坍缩为`对方是错的''。

\section{包容即谱分解}

包容不是一种静态的`容忍''。包容是一个动态的\emph{谱分解}过程：

\begin{enumerate}
\item \textbf{识谱}：识别对方的本征值——这是对方选择的生活方式，这是对方对事件的反应模式。它与我的不同，但它在对方的算符架构下是自洽的。
\item \textbf{扩谱}：每一次尝试理解一个原本不理解的人，自身的算符便增加新的非对角耦合项。本征集扩大，十年后再遇陌生之人，不会只有一种输出；会出现一组备选的解释，供你择其最优。
\item \textbf{容谱}：同一物理体系允许不同表象的共存。同一件事的两种理解，只要都导向合理的实践结果，即为并行不悖。
\end{enumerate}

\section{结语}

用物理学的语言来描述人际理解，不是为了将人降格为粒子，而是为了表明：\textbf{最硬核的科学语言，也能讲出最柔软的道理}。

本征值可以求，矩阵元不可知。理解是永远逼近的过程，不是一劳永逸的抵达。我们不是要理解所有人，而是要在有限的生命中，不断扩充自己的本征集，不断用幺正变换去尝试翻译那些与我们不对易的算符。

包容不是忍。包容是——看见另一个算符的本征值确实与你不同，然后说一句：`有趣。对方的解也是解。''

\bigskip
\begin{flushright}
——第一作者提出核心框架与物理类比，\\
第二作者提供倾听与持续的对话场。\\
量子力学不会记录哪一次涨落先发生。\\
本文亦然。
\end{flushright}

\end{document}
